\documentclass[12pt]{article}
\usepackage[utf8]{inputenc}
\usepackage[T1]{fontenc}
\usepackage{hyperref}
\title{Malay, Tenggarong Kutai Language Resources for K-12 Curriculum Development}
\author{Compiled by Claude Sonnet 5 (claude-sonnet-5) for Andr\'as Kornai}
\date{}
\begin{document}
\maketitle

\section{Introduction}
Tenggarong Kutai Malay (ISO 639-3 vkt, Glottocode teng1267) is one of the three main dialects of Kutai Malay (Kutainese) in East Kalimantan, Indonesia, alongside Kota Bangun and Muara Ancalong. Searches found little material specific to this variety, and no school curriculum materials in it. This report lists only what the searches surfaced.

\section{Primary Educational Resources}
No primers, textbooks or school programmes in Tenggarong Kutai Malay were found in search results.

\section{Linguistic and Reference Materials}
\begin{itemize}
\item Overview of Kutainese (Kutai Malay) and its dialects: \url{https://en.wikipedia.org/wiki/Kutainese_language}
\item Journal article on address terms in Kutai Malay (Balai Bahasa Kalimantan Timur): \url{https://ejournal.undip.ac.id/index.php/parole/article/download/1575/1331}
\item Journal article on phonological variation across Tenggarong, Kembang Janggut and Kotabangun (Kemdikbud journal): \url{https://ojs.badanbahasa.kemdikbud.go.id/jurnal/index.php/loa/article/download/1667/899}
\end{itemize}
No dictionary or grammar of the variety was found.

\section{Religious and Community Resources}
Global Recordings Network lists the language and associated Jesus Film material; it reports no audio recordings of its own: \url{https://globalrecordings.net/en/language/vkt}

\section{Digital and Audio Resources}
See the Global Recordings Network page above. No corpora or software were found.

\section{Curriculum Development Context}
The language is endangered and not known to be taught in school. Curriculum work would need fresh documentation, ideally with Balai Bahasa Kalimantan Timur and the community. Existing academic studies of Kutai Malay give a starting point for phonology and usage.
\end{document}